\guidepage{The model: boundary restoration and Boolean degree}
A state has logical input wires $x$, zero-initialized ancillary wires, and a
target bit $y$. A restored oracle is a legal classical circuit satisfying
\[
 C(x,0^a,y)=(x,0^a,y+f(x))
 \qquad\text{for every }x\text{ and both }y.
\]
The input and scratch restoration conditions hold at the boundary. Wires may
change internally, and the target may control an intermediate gate. Each
primitive gate is an involution:
\[
 y\leftarrow y+1,\quad y\leftarrow y+x_i,\quad
 y\leftarrow y+x_i x_j.
\]
In a CNOT the control differs from the target. In a Toffoli all three wires
are distinct. These are construction rules, not informal assumptions.

All polynomial arithmetic in the Boolean state model is over $\F_2$, with
$x_i^2=x_i$. Every function has a unique multilinear algebraic normal form.
Its degree is the largest support size with a nonzero coefficient; the zero
function has mathematical degree $-1$ in the Python helpers. A formal degree
bound uses a natural number, so zero has a degree-zero certificate.

\begin{tabularx}{\textwidth}{@{}p{1.15in}X@{}}\toprule
Quantity & Meaning\\\midrule
$T$ & Number of Toffoli gates; NOT and CNOT have zero cost in this quantity\\
$L$ & Total number of NOT, CNOT, and Toffoli gates\\
$M$ & Number of binary AND operations in an XOR--AND computation\\
$n$ & Actual number of input bits, including all variables in a displayed family\\
$\dim_K A$ & Dimension of a specified linear summary algebra over $K$\\\bottomrule
\end{tabularx}

A bound on one row of this table cannot silently be substituted for a bound
on another row. In particular, a computation can have few Toffolis and many
CNOTs. An exponentially large ambient algebra can have succinct expressions.
Those distinctions explain the projection and gluing examples later.

\sourcepath{docs/FOUNDATIONS.md, Conventions; formalization/lean/CircuitObstructions/Gates.lean}

\guidepage{An arbitrary cut compares two reversible computations}
Write the circuit as $C=Q\circ P$. The prefix $P$ runs forward; the inverse
suffix $Q^{-1}$ starts from the same zero-target boundary state. Define
\[
 A(x)=P(x,0),\qquad B(x)=Q^{-1}(x,0),\qquad
 D(x)=P(x,1)+P(x,0).
\]
The zero ancillas are suppressed here. The two executions generally arrive
at different interior states, even though they start at matching boundaries.

\visualfigure{cut_transport}{The exact C01 objects. Both arrows remain inside the legal reversible state space.}

Correctness on both target branches gives
\[
 Q^{-1}(x,b)=P(x,b+f(x)).
\]
To derive it, apply $C$ to $(x,b+f(x))$ and then apply $Q^{-1}$ to the result.
For $b=0$, $B=A$ when $f(x)=0$, while $B=P(x,1)$ when $f(x)=1$. Thus
\[
 H(x):=A(x)+B(x)=f(x)D(x).
\]
The two original states differ in their target. Since $P$ is injective,
$D(x)\ne0$. Each coordinate discrepancy is supported on positives, and
some coordinate discrepancy detects every positive input.

\evidence{General written identity; archived formal declarations
\code{cut_transport_identity}, \code{direction_nonzero}, and \code{complete_safe_cover}.}

\guidepage{Why completeness is pointwise, and why no decoder is needed}
For a coordinate $i$, let $H_i=A_i+B_i$. The cut identity says
\[
 H_i(x)=f(x)D_i(x),\qquad H_i\le f,\qquad \bigvee_i H_i=f.
\]
The last equation is pointwise. For every fixed accepted $x$, at least one
coordinate has $D_i(x)=1$. The detecting coordinate may depend on $x$ and
may be an input, scratch, or target wire. The argument does not require a
single fixed coordinate to detect all positives.

At an empty prefix, the direction is simply the target unit vector. At a
full prefix, all non-target wires have restored and the target discrepancy
is $f$. Interior cuts can spread the direction through the workspace.
These minimum cases verify the composition order and the role of the two
target branches.

The zero-input cases are also valid. With only a target wire, the empty
circuit realizes the constant-zero oracle and one NOT realizes the
constant-one oracle. At every cut, the direction is nonzero and the
discrepancy cover has the correct accepting set.

\begin{quote}\small
\textbf{Essential hypothesis.} The endpoint condition holds for both initial
target values. The inverse-suffix derivation uses the branch $b+f(x)$.
Correctness only for initial target zero is not the stated model.
\end{quote}

The theorem also does not require scratch to remain zero at the cut. Imposing
that extra condition would exclude admitted restored circuits. Conversely,
dropping target separation from the boundary encoding permits a vacuous
endpoint equation: an encoding that ignores $y$ can make the empty circuit
appear to compute a nonzero function. The formal \code{Boundary.target_eq}
property prevents precisely this degeneracy.

\sourcepath{CutComparison.lean: inverse_suffix_boundary, every_index_cut, canonical_complete_safe_cover}

\guidepage{A dirty scratch wire can be fully restored}
Consider the four-wire oracle in the Lean examples. First compute $ab$ into
scratch $s$, copy $s$ into the target, and uncompute $s$:
\[
 (a,b,s,y)\xrightarrow{T(a,b;s)}(a,b,s+ab,y)
 \xrightarrow{C(s;y)}(a,b,s+ab,y+s+ab)
 \xrightarrow{T(a,b;s)}(a,b,s,y+s+ab).
\]
The displayed final expression specializes to the required oracle when
the boundary scratch is zero. On the positive input $a=b=1$, the scratch
is dirty between the first and last gates.

\visualfigure{dirty_and_trace}{Exact states for the accepting input, in wire order $(a,b,s,y)$. White cells are zero; teal cells are one.}

The two panels check both initial targets. The final target is toggled,
the logical inputs are unchanged, and scratch returns to zero. A proof that
used an intermediate clean-workspace assumption would fail to cover this
explicit legal fixture.

This circuit uses two Toffolis. It is a useful restoration witness, not the
optimal two-input AND oracle: the direct gate $T(a,b;y)$ computes that oracle
with one Toffoli and no scratch. The guide distinguishes an exact example
from an optimization claim.

\evidence{Exact new trace assertion imported from the source gate semantics;
the same fixture is formalized by \code{Examples.andOracle_restores}.}

\guidepage{The output direction can move away from the target}
Two data/target wires already suffice. A SWAP is implemented by three legal
CNOTs. Let $P$ be that SWAP, and let $Q$ be NOT on the data wire followed
by another SWAP. Then
\[
 (x,y)\xrightarrow{P}(y,x)
 \xrightarrow{Q}(x,1+y).
\]
The whole circuit is a restored constant-one oracle with zero Toffolis.
At the cut after the first SWAP,
\[
 A=P(x,0)=(0,x),\qquad B=Q^{-1}(x,0)=(1,x),
 \qquad H=(1,0).
\]
The target discrepancy is zero on every input, while the data discrepancy
is one on every input. Reading only the target at this cut would reject
the constant-one function.

\begin{center}
\begin{tabular}{@{}ccccc@{}}\toprule
$x$ & $A$ & $B$ & $H$ & named target discrepancy\\\midrule
0 & $(0,0)$ & $(1,0)$ & $(1,0)$ & 0\\
1 & $(0,1)$ & $(1,1)$ & $(1,0)$ & 0\\\bottomrule
\end{tabular}
\end{center}

This example uses the target as a control during each SWAP. That behavior
is expressly allowed by the model. Its relevance is logical: completeness
belongs to the full register comparison, not to an arbitrarily chosen wire.
The coordinate packets remain perfectly safe even though the detection
direction has moved.

\evidence{Constructive formal counterexample to target-only reconstruction:
\code{outputActiveOracle_restores}, \code{moved_discrepancy_data},
\code{moved_discrepancy_target}.}

\guidepage{Degree bounds must retain shared register products}
Take four independent input bits. Compute
\[
 q_1=x_1x_2,\qquad q_2=x_3x_4,\qquad q_3=q_1q_2.
\]
After two ANDs, each nonlinear register has degree two, but their product
has degree four. After the third AND, a coordinate of degree four appears.

\visualfigure{sharing_degree}{Exact degrees from the complete four-input truth tables. The square marker is a two-register product, not a coordinate.}

The correct invariant tracks a product of $r$ current registers after $t$
Toffolis:
\[
 \deg\!\left(\prod_{i\in L}s_i(x)\right)\le t+|L|=t+r.
\]
At $t=2$, the product $q_1q_2$ has $r=2$ and saturates $t+r=4$. At $t=3$,
the resulting coordinate has $r=1$ and saturates $t+1=4$.

This explains why tracking only individual register degrees and then
adding their separate histories can lose the shared accounting. The
invariant applies to the actual evolving state. It does not unfold a
reused register into a new copy of the entire computation.

\evidence{Exact finite fixture plus the general register-product invariant;
the source proves \code{workspaceDegree_run} for arbitrary gate lists.}

\guidepage{The invariant handles duplicates, NOT, and CNOT}
Suppose a gate toggles target $u$ by a product of $c$ controls. In a product
list $L$ that contains $u$, remove all occurrences of $u$ to obtain $R$.
Boolean idempotence leaves one updated target factor, and
\[
 \prod_{i\in L}s_i'=
 s_u\prod_{i\in R}s_i+
 \left(\prod_{j\in\mathrm{controls}}s_j\right)\prod_{i\in R}s_i.
\]
The old invariant bounds these terms by $t+1+|R|$ and $t+c+|R|$.
Since $|R|+1\le |L|$, the budget rises by at most $\max\{c-1,0\}$.

\begin{center}
\begin{tabular}{@{}lccc@{}}\toprule
Gate & $c$ & toggle expression & extra budget\\\midrule
NOT & 0 & 1 & 0\\
CNOT & 1 & one control & 0\\
Toffoli & 2 & two controls & 1\\\bottomrule
\end{tabular}
\end{center}

If the target is absent from $L$, the product is unchanged. An empty list
is the constant-one product. Repeated target indices are removed using
idempotence; repeated other indices can only reduce multilinear degree.
Controls already appearing in $R$ also cause no difficulty.

For general XOR--AND computations, the written proof introduces a Boolean
variable for each AND output. Every wire is affine in inputs and preceding
gate variables. Substitute those variables in reverse topological order.
A squarefree monomial contains each substituted variable at most once,
so replacing it by a quadratic raises degree by at most one. After $t$
substitutions, the affine output has degree at most $t+1$.

These are two precise proof routes. The general syntax argument is written;
the maintained NCT register argument is in Lean. The formalization does not
assert a separate arbitrary XOR--AND compiler or a canonical ANF algorithm.

\sourcepath{MultiplicativeDegree.lean; docs/FOUNDATIONS.md, F3}

\guidepage{One safe packet need not cover every accepting point}
For $k=3$ consider
\[
 F_3(a,b,z_1,z_2,z_3)=ab\ \mathrm{OR}\ z_1z_2z_3.
\]
The packet $ab$ is safe and has degree two. It detects every teal point
below. It misses three additional positives, which require the other
packet $z_1z_2z_3$.

\visualfigure{packet_coverage}{All 32 assignments. There are 11 positives: 8 in the teal row and 3 additional gold points. The coordinate labels use $(z_1,z_2,z_3)$ order.}

The one-packet parameter $\alpha_1$ asks for any nonzero safe polynomial.
The complete parameter $\beta$ asks for safe detection of every positive:
\[
 \beta_v(f)=\min\{\deg h:h\le f,\ h(v)=1\},\qquad
 \beta(f)=\max_{v:f(v)=1}\beta_v(f).
\]
For this monotone function, its minimal positive supports have sizes two
and three. Thus $\alpha_1(F_3)=2$ while $\beta(F_3)=3$. Selecting one
easy accepting region loses the pointwise obstruction at the other region.

\evidence{All 32 points and the two parameter values were freshly recomputed;
the general monotone minterm characterization is a written proof for nonzero functions.}

\guidepage{Balanced cuts turn coverage into a Toffoli bound}
Choose a cut after $\lfloor T/2\rfloor$ Toffolis. Both sides have at most
$\lceil T/2\rceil$ Toffolis. Reversing the suffix preserves its count because
every gate is self-inverse. Each forward or inverse-suffix coordinate has
degree at most $\lceil T/2\rceil+1$, and their XOR has the same upper bound.

For every positive $v$, completeness chooses a discrepancy $H_i$ that
detects $v$. Its safety and degree bound imply
\[
 \beta(f)\le\left\lceil\frac T2\right\rceil+1,
 \qquad 2\beta(f)\le T+3.
\]
The integer rounding matters. If $T=2r$, the first bound gives
$2\beta\le T+2$. If $T=2r+1$, it gives $2\beta\le T+3$.
The uniform bound uses the latter constant.

An affine output mask $\ell(x)$ costs zero additional Toffolis: append
CNOTs from restored inputs and, if needed, a target NOT. This constructs
a restored oracle for $f+\ell$, so the bound applies to every affine mask.
Define
\[
 \beta^*(f)=\max_{\ell\text{ affine}}\beta(f+\ell).
\]
Together with the output degree bound,
\[
 \boxed{T\ge\max\{0,\deg f-1,2\beta^*(f)-3\}.}
\]

There are finitely many affine masks. The Lean definition uses the least
upper bound of their beta values, which equals this finite maximum. Its
affine-output theorem constructs legal gates rather than assuming masks
are free transformations of an arbitrary interior computation.

\evidence{Written theorem and archived endpoint
\code{restored_degree_and_cover_bound}; the dual is a separate written result.}

\guidepage{An affine shift reveals a harder accepting point}
For $k\ge2$, let $B=z_1\cdots z_k$ and $F_k=ab\ \mathrm{OR}\ B$ on
$n=k+2$ inputs. Add the affine mask $a$ to obtain $G_k=F_k+a$.
On the face $b=1$,
\[
 G_k|_{b=1}=(1+a)B.
\]
This is a singleton indicator on the $k+1$ variables $(a,z)$.
At $v=(a=0,b=1,z=1^k)$, every safe detecting polynomial must restrict to
that singleton indicator, so its degree is at least $k+1$.

The general two-positive indicator construction gives the matching upper
bound $\beta(G_k)\le n-1=k+1$. The maintained family proves
\[
 \max_\ell\alpha_1(F_k+\ell)=3,\qquad \beta^*(F_k)=k+1.
\]
The gap between one-packet and complete degree therefore grows with $k$.

\begin{center}
\begin{tabular}{@{}rrrrr@{}}\toprule
$k$ & $n$ & $\beta^*$ & complete bound $2k-1$ & combined degree/one-packet bound\\\midrule
2 & 4 & 3 & 3 & 3\\
3 & 5 & 4 & 5 & 4\\
4 & 6 & 5 & 7 & 5\\
6 & 8 & 7 & 11 & 7\\\bottomrule
\end{tabular}
\end{center}

For $k=3$, all 64 affine masks were freshly checked and the optimized
parameters are exactly $(3,4)$. The all-size family proof is not inferred
from those 64 cases. Its restored construction has $T\le2n-3$, so this
remains a linear-size easy family. The exact optimum between the lower and
upper constructions is not asserted when they differ.

\sourcepath{docs/THEORY.md, section 2; reviews/visual_examples.json, unbalanced_or}

\guidepage{A degree-one dual certificate for two-input AND}
Let $f(a,b)=ab$, with unique positive point $(1,1)$. For degree at most one,
the evaluation vector is $v_1(a,b)=(1,a,b)$. The three negative points
already span the positive vector.

\visualfigure{evaluation_dual}{The first three rows are negative inputs. Their XOR equals the last evaluation row.}

Any affine polynomial is a coefficient covector $c=(c_0,c_a,c_b)$.
Safety means that it pairs to zero with each negative row. The displayed
dependence therefore forces its value at $(1,1)$ to be zero. No safe affine
packet detects the positive point.

The degree-two polynomial $ab$ does detect it, so
$\beta_{11}(f)=\beta(f)=2$. Since $\beta^*\le n=2$ and $\beta^*\ge\beta$,
$\beta^*=2$. The complete bound gives $T\ge1$, attained by one direct
Toffoli to the target.

The constant coordinate also checks parity: this dependence uses three
negative rows, an odd number. Without the constant coordinate, that parity
conclusion would not follow.

\evidence{Exact finite dependence plus written duality. Lean proves the
corresponding AND2 degree and beta exclusions by four-corner parity,
without formalizing the general evaluation-span dual.}

\guidepage{The finite evaluation-span dual in full}
Fix a positive $v$ and integer $d\ge0$. Let
\[
 v_d(x)=(x^S)_{S\subseteq[n],\ |S|\le d},\qquad
 U=f^{-1}(0),\qquad N=\Span_{\F_2}\{v_d(u):u\in U\}.
\]
The constant monomial $S=\varnothing$ is included. Degree-$d$ coefficient
covectors form the finite dual of this evaluation space. A covector is safe
exactly when it annihilates $N$.

If $v_d(v)\in N$, every safe covector vanishes at $v$. If $v_d(v)\notin N$,
extend a basis of $N$ by $v_d(v)$ and then to a basis of the full space.
Assign zero to the $N$ basis and one to $v_d(v)$. The resulting covector is
a safe detector. Hence
\[
 \boxed{\beta_v(f)>d\quad\Longleftrightarrow\quad v_d(v)\in N.}
\]
Choose a basis of $N$ from negative evaluation rows. A span representation
of the positive row gives a subset $W\subseteq U$ with
\[
 v_d(v)=\sum_{u\in W}v_d(u),\quad |W|\text{ odd},\quad
 |W|\le\sum_{j=0}^{\min(d,n)}\binom nj.
\]
It is a replayable finite obstruction to every safe degree-$d$ detector.

At $d=0$, a negative exists exactly when the constant-one detector is unsafe.
At $d\ge n$, the full evaluation rows are independent and no positive row
lies in the negative span. This matches $\beta_v\le n$ from point indicators.
For $f=0$, there is no positive $v$ and the separate convention is $\beta(0)=0$.

Fresh enumeration checked all 278 functions on zero through three inputs,
all positives, and every $d=0,\ldots,n$: 4,201 primal/dual checks agree.
Those bounded checks support the independently supplied universal proof.

\sourcepath{scripts/core_dual_checks.py; reviews/core_dual_checks.json}

\guidepage{Every function through three inputs has a matching oracle}
The small calibration concerns every truth table on one, two, or three
inputs. It is not a search through all reversible circuit words. A theorem
lower bound and an explicit construction coincide.

\visualfigure{small_oracle_counts}{Reconstructed source constructions for all 276 functions. Totals are 28 with $T=0$, 120 with $T=1$, and 128 with $T=3$.}

Affine functions use target CNOTs and possibly NOT. A nonaffine quadratic
through three bits has alternating quadratic rank two, so it is a product
of two independent linear forms plus an affine residual. Compute those
forms on clean wires, apply one Toffoli, uncompute, and append the residual.

Every three-bit cubic is a translated three-variable monomial plus an
affine residual. The singleton monomial needs three Toffolis by complete
coverage, and compute/toggle/uncompute attains that number.

The new figure script rebuilds and validates all 276 circuits and both target
branches, covering 4,240 scalar input/target cases via exact truth tables.
Its totals match the frozen baseline. Zero-input constants were checked in
the separate dual suite, rather than being added to this 276-function count.

\evidence{Matching written lower bounds and exact constructed implementations;
not exhaustive optimal-circuit synthesis. The generic monomial cost theorem
is not newly Lean-certified by this edition.}

\guidepage{Projection is defined in the original coordinates}
For $m\ge2$, put $a=(1,\ldots,1)$ and
\[
 V=\F_2^m\setminus\{0,a\},\qquad n=2^m-2.
\]
There is one input bit $x_v$ for each $v\in V$. Define
\[
 L_j(x)=\sum_{v:v_j=1}x_v,\qquad
 P_m(x)=\prod_{j=1}^m L_j(x).
\]
The linear map $x\mapsto L(x)$ has rank $m$, because the standard basis
columns are present. In the $L$ coordinates, $P_m$ is a monomial in $m$
independent bits. Its restored nonlinear cost is $2m-3$.

Now extract the quadratic homogeneous ANF component $Q_m$ in the
\emph{original} $x_v$ coordinates. Since $P_m(0)=P_m(e_v)=0$, for distinct
$u,v$ its quadratic coefficient is
\[
 [x_ux_v]P_m=P_m(e_u+e_v)=[u+v=a].
\]
The involution $v\mapsto v+a$ pairs every variable with one other variable:
\[
 Q_m=\sum_{\{v,v+a\}}x_vx_{v+a}.
\]
It is a quadratic inner product on $n/2$ disjoint pairs. Its alternating
polar matrix has full rank $n$.

The coordinate choice is essential. Homogeneous grade is not invariant
under arbitrary affine changes of variables. A cheap representation in
the $L$ coordinates does not imply cheap independent generation of the
quadratic grade in the original coordinates. The comparison keeps the
definition of projection fixed throughout.

\sourcepath{docs/THEORY.md, section 3; checks/verify_projection.py}

\guidepage{Exact nonlinear projection costs versus actual input length}
The source proves the all-size formulas
\[
 M(P_m)=m-1,\quad T_{\mathrm{rest}}(P_m)=2m-3,\qquad
 M(Q_m)=T_{\mathrm{rest}}(Q_m)=n/2.
\]
The chart uses actual input length $n=2^m-2$ on its horizontal axis.

\visualfigure{projection_costs}{Exact formulas for $m=2,\ldots,10$. Both axes are logarithmic; the plotted points are not optimal-circuit search results.}

For $Q_m$, each disjoint quadratic pair is accumulated directly into the
target with one Toffoli. The lower bound allows nonlinear intermediates:
any computation with $t$ multiplications is affine on an affine subspace
of codimension at most $t$. A nondegenerate quadratic can be affine only
on a subspace with isotropic direction space, whose dimension is at most
$n/2$. Thus $t\ge n/2$.

The separation grows exponentially in $m$, but $n$ itself grows
exponentially in $m$. The quadratic cost is linear in actual input length.
The conclusion is about nonlinear generation and a specified grade
extraction operation, not a superpolynomial unrestricted circuit bound.

\evidence{All-size written cost proofs; constructor gate counts freshly
asserted through $m=10$. No full truth-table enumeration at $m>4$ is claimed.}

\guidepage{At $m=10$, the total-gate accounting points the other way}
Here $n=1,022$. The full function uses 17 Toffolis; its quadratic grade
needs 511. However, computing and uncomputing the $m$ parities also uses
\[
 2m(2^{m-1}-1)=20\cdot511=10,220\quad\text{CNOTs}.
\]
The explicit full construction has $10,237$ total gates. The quadratic
accumulation uses only 511 total gates.

\visualfigure{projection_m10}{Two different resource comparisons at the same input length. The right panel reports the totals of the stated constructions, not an asserted optimal total for $P_{10}$.}

The left panel is a valid nonlinear-cost separation. The right panel
prevents interpreting it as an exponential ordinary total-gate separation.
The total full construction is $O(n\log n)$, whereas the quadratic
construction is $n/2$ gates.

A proposed invariant that charges an original computation for every
separately extracted grade would have to justify that extraction's cost.
These formulas refute an overhead bounded only by a polynomial in the
original nonlinear count, independently of input and linear resources.
They do not exclude an overhead polynomial in $n+T$ or a different
precisely defined invariant.

\evidence{Exact counts obtained from the source constructors; optimal
Toffoli formulas have written proofs. The full total-gate count is an upper
construction count.}

\guidepage{Quadratic projection can destroy safety}
At $m=3$, variables are ordered by source labels $v=1,2,3,4,5,6$.
Assignment index $13$ is read least-significant-bit first:
\[
 (x_1,x_2,x_3,x_4,x_5,x_6)=(1,0,1,1,0,0).
\]
The active labels have syndrome $1+3+4=6$, not $a=7$.

\visualfigure{unsafe_projection}{A concrete accepted point for $Q_3$ that is rejected by $P_3$. The variable ordering matches the source checker.}

The three linear forms evaluate to $L=(0,1,1)$, so $P_3=L_1L_2L_3=0$.
The quadratic grade has pairs $(1,6)$, $(2,5)$, $(3,4)$, so
\[
 Q_3=x_1x_6+x_2x_5+x_3x_4=0+0+1=1.
\]
Thus $Q_3\notin I_{P_3}$: homogeneous extraction does not preserve the
positive-support ideal. The whole polynomial is supported on its own
accepting set, but its isolated homogeneous components need not be.

The new script checks all 64 inputs here, finds 24 false positives, and
confirms that 13 is the first in the stated integer ordering. The frozen
$m=4$ control similarly has 7,616 false positives among its $2^{14}$
assignments. These are separate finite statements with exact ranges.

\evidence{Fresh complete $m=3$ truth-table check; baseline $m=4$ check was
reproduced by the full source verifier. Safety failure is already established
by the single valid assignment above.}

\guidepage{Witness indicators are orthogonal idempotents}
Let $K$ be a field and $W=\{0,1\}^m$. The full witness algebra is
$B=K^W$ with pointwise addition and multiplication. For each witness $v$,
the indicator $\delta_v$ is one at $v$ and zero elsewhere. Then
\[
 \delta_v^2=\delta_v,\qquad \delta_u\delta_v=0\ (u\ne v),
 \qquad \sum_v\delta_v=1.
\]

\visualfigure{orthogonal_gluing}{The four-witness multiplication table at $m=2$. Each nonzero diagonal entry is its own point indicator.}

The four indicator functions form a basis of $K^4$. General $m$ has $2^m$
such basis elements. A verifier is the sum of the indicators of its accepted
witnesses, and pointwise multiplication models intersection of accepting
sets.

This algebra retains exact witness compatibility: disjoint singleton
predicates multiply to zero. A summary that preserves that product operation
while retaining every singleton must preserve the same orthogonal structure.
The following dimension theorem says exactly how much fixed linear space
that requires.

\evidence{Exact $m=2$ multiplication table; general product-algebra identities
hold over any field. No characteristic-two assumption is needed for the
dimension proof.}

\guidepage{Faithful linear multiplicative summaries need $2^m$ dimensions}
Suppose $\Phi:K^W\to A$ is $K$-linear and multiplicative, and
$\Phi(\delta_v)\ne0$ for every $v\in W$. Put $e_v=\Phi(\delta_v)$.
The images are nonzero orthogonal idempotents.

If $\sum_v c_v e_v=0$, multiply by $e_u$. Orthogonality removes every other
term and idempotence gives $c_u e_u=0$. Since $K$ is a field and $e_u\ne0$,
$c_u=0$. Hence all $e_v$ are linearly independent and
\[
 \boxed{\dim_K A\ge |W|=2^m.}
\]
No unital hypothesis is needed. At $m=0$, there is one empty witness and
the bound is simply $\dim_K A\ge1$.

For a unital scalar map to $K$, scalar idempotents are zero or one.
Orthogonality permits at most one nonzero image, and preservation of one
requires exactly one. Such a map is evaluation at one witness. For $m\ge1$
it loses at least one other singleton verifier. For $m=0$, there is no other
witness to lose.

\begin{quote}\small
\textbf{Model boundary.} The theorem quantifies a fixed summary serving all
predicates at a witness length, preserving linear combinations and exact
pointwise multiplication. An unrestricted algorithm need not use such a
summary. Syntax, nonlinear encodings, changing spaces, context-dependent
operations, and succinct expressions are outside these hypotheses.
\end{quote}

An exponentially large algebra can still contain short circuit expressions.
Therefore exponential dimension does not imply exponential expression length
or a lower bound for every algorithm on one explicit decision family.

\sourcepath{docs/THEORY.md, Theorem 4.1; reviews/extended_audit.md}

\guidepage{Distinct witness monomials prevent formal cancellation}
Over a characteristic-two field, define
\[
 P_R(z)=\sum_{w\in\{0,1\}^m}R(w)\prod_i z_i^{w_i}.
\]
The $z_i$ are formal indeterminates, not Boolean variables. Different
witnesses produce different monomials, so $P_R=0$ exactly when there is
no accepting witness.

\visualfigure{witness_fingerprint}{An exact three-bit verifier accepts witness indices $1,2,5,6$ in least-significant-bit order. Four accepts have parity zero, but their monomials are distinct.}

Here
\[
 P_R=z_1+z_2+z_1z_3+z_2z_3=(z_1+z_2)(1+z_3)\ne0.
\]
The ordinary unweighted characteristic-two sum of the four accepts is
zero. The formal polynomial avoids this cancellation by retaining witness
labels as monomial supports.

Nonzero as a formal polynomial does not mean nonzero at every field point.
For instance, this example vanishes whenever $z_1=z_2$ or $z_3=1$.
Random evaluation is governed by a root bound, not by the formal identity
alone. Evaluating the polynomial efficiently remains a separate task.

\evidence{Exact support fixture; general no-cancellation argument is written.
The source finite evaluator retains $2^m$ coefficients.}

\guidepage{A short expression inside a trace is not a computed trace}
Let $\tau_w$ extract the top multilinear coefficient in the Boolean witness
ring. For any $G(w)$ over a characteristic-two field,
\[
 \tau_w(G)=\sum_{v\in\{0,1\}^m}G(v).
\]
A proper-support monomial evaluates to one an even number of times; the
full-support monomial does so once. This proves the identity directly.

Applying it to a weighted expression gives
\[
 P_R(z)=\tau_w\!\left[
 R(w)\prod_{i=1}^m\bigl(1+(1+z_i)w_i\bigr)\right].
\]
At $w_i=0$, a multiplier equals one. At $w_i=1$, it equals $z_i$.
Summing evaluations therefore produces exactly the witness monomials.
For the preceding verifier, the four surviving products are
$z_1,z_2,z_1z_3,z_2z_3$.

The circuit expression \emph{inside} the trace adds only $O(m)$ operations
to the verifier. That observation bounds the expression's construction,
not the cost of applying $\tau_w$. The finite checker computes with a
complete table of $2^m$ witness coefficients.

\begin{center}
\begin{tabularx}{\textwidth}{@{}p{1.8in}X@{}}\toprule
Established & Not supplied by the expression\\\midrule
Small weighted circuit inside $\tau$ & A polynomial-time evaluator for the trace\\
Formal polynomial detects nonemptiness & A deterministic identity test from one evaluation\\
Exact finite table algorithm & Polynomial storage in the witness length\\\bottomrule
\end{tabularx}
\end{center}

At $m=0$, the empty product is one, the top coefficient is the constant
itself, and there is a single empty witness. The definitions remain valid.

\sourcepath{docs/THEORY.md, equations (10)--(11); checks/verify_fingerprint.py}

\guidepage{Disjoint-variable factorization uses each predicate's own block}
Suppose $A$ depends on witness variables in block $S$ and $B$ on a disjoint
block $T$. Write $\tau_{z_S}(A)=\sum_s A(s)z_S^s$ for the weighted trace
on $S$, distinct from the unweighted top-coefficient notation $\tau_w$.
Weighted sum over the union factors as
\[
 \sum_{s,t} A(s)B(t)z_S^s z_T^t
 =\left(\sum_s A(s)z_S^s\right)
  \left(\sum_t B(t)z_T^t\right).
\]
Each factor is traced on its own variable block. The union trace is then
their product. This is an ordinary factorization of a finite double sum.

For $A=w_1$ and $B=w_2$,
\[
 \tau_{z_1}(A)=z_1,\qquad \tau_{z_2}(B)=z_2,\qquad
 \tau_{z_1,z_2}(AB)=z_1z_2.
\]
Tracing both factors over the full two-bit ambient domain introduces unused
variable sums:
\[
 \tau_{z_1,z_2}(A)=z_1(1+z_2),\qquad
 \tau_{z_1,z_2}(B)=z_2(1+z_1).
\]
Their product generally differs from the union trace. At $z_1=z_2=1$, each
full ambient factor is zero, while the trace of $AB$ is one. Disjoint
dependence alone does not justify multiplying traces with mismatched domains.

The source checker tests all four one-bit predicates per block and every
pair of GF(16) weights: $4\cdot4\cdot16^2=4,096$ own-block factorization
cases. It does not claim an ambient-domain product identity.

\evidence{General written factorization under the own-block convention;
exact finite source check and an explicit ambient-domain counterexample.}

\guidepage{Overlap breaks the attempted context-free product rule}
One shared witness bit already defeats multiplicativity. With
$\tau_z(P)=P(0)+zP(1)$,
\[
 \tau_z(w)=z,\qquad \tau_z(1+w)=1,
 \qquad \tau_z(w(1+w))=0.
\]
The product of the separate summaries is $z$. Also $w^2=w$ in the Boolean
ring, so $\tau_z(w^2)=z$, whereas $\tau_z(w)^2=z^2$.

\visualfigure{overlap_gf16}{A concrete field point $z=\theta$ with $\theta\ne0,1$. The two columns are unequal in both rows.}

The field is $\F_{16}=\F_2[\theta]/(\theta^4+\theta+1)$. The checker stores
field elements as four-bit coefficient vectors: $\theta$ is integer 2 and
$\theta^2$ is integer 4. These integers encode field elements; they are not
ordinary integer multiplication. Exact field multiplication gives the
reported mismatch.

The first row loses witness compatibility: $w$ and $1+w$ have no common
accepting witness, yet their separate weighted summaries multiply to a
nonzero value. The second row loses Boolean idempotence. This is a genuine
counterexample to a context-free multiplicative trace compiler, not a
floating-point effect or an implementation anomaly.

\evidence{Exact GF(16) fixture imported from the source field routines. The
general overlap failure is symbolic and does not depend on the selected field.}

\guidepage{Randomized detection is conditional on efficient evaluation}
The witness polynomial has total degree at most $m$. Over a finite field
$K$, a nonzero polynomial of total degree at most $m$ vanishes on at most
an $m/|K|$ fraction of independent uniform field points. The standard
inductive proof separates the leading coefficient in the last variable
and applies the univariate root bound.

If a uniform polynomial-time exact evaluator for the weighted trace were
available, with efficient field representation and arithmetic and
$|K|\ge4m$, one evaluation would accept a satisfiable verifier with
probability at least $3/4$. An unsatisfiable verifier has the zero polynomial,
so it never produces a false acceptance. Repetition with independent samples
reduces the false-rejection probability to at most $4^{-r}$ after $r$ trials.

\begin{center}
\begin{tabular}{@{}rrrr@{}}\toprule
Witness bits & GF(16) points & maximum zeros in source tests & degree-bound zeros\\\midrule
1 & 16 & 1 & 1\\
2 & 256 & 31 & 32\\
3 & 4,096 & 721 & 768\\\bottomrule
\end{tabular}
\end{center}

These source tests enumerate every verifier at each displayed witness length
and all field points for the polynomial evaluations. The independent
coefficient-extraction comparison has a narrower three-variable grid:
coordinates $0,1,2,3$, rather than all GF(16) points. Both scopes are recorded.

This is a conditional randomized algorithmic implication. The required
efficient trace evaluator is not supplied. The existing exact evaluator is
exponential in witness length, and no deterministic $P=NP$ conclusion is
drawn. Random bits and polynomial output description alone do not account
for the evaluation operation.

\sourcepath{docs/THEORY.md, section 5; evidence/baseline/fingerprint.json}

\guidepage{Reproduce the figures and inspect their exact evidence}
All 12 figure pairs are vector PDFs plus raster PNG previews. Their single
generation command, from the publication-set root, is
\begin{lstlisting}
python scripts/make_figures.py
\end{lstlisting}
The script imports the original gate, profile, construction, and GF(16)
functions without invoking their release-writing entry points. It writes
only new figures and \code{reviews/visual_examples.json}. That record contains
fixture truth-table integers, ordered bits, exact gate lists, scalar counts,
the source hashes, and the distinction between finite assertions and
plotted all-size formulas.

The separate dual sanity suite runs as
\begin{lstlisting}
python scripts/core_dual_checks.py \
  --output reviews/core_dual_checks.json
\end{lstlisting}
It compares exhaustive degree-bounded safe truth tables with a separate
Gaussian-elimination calculation of the negative evaluation span. Its
range is all truth tables at $n=0,1,2,3$, not an extrapolation.

\begin{tabularx}{\textwidth}{@{}p{1.4in}X@{}}\toprule
Evidence layer & Publication treatment\\\midrule
General mathematics & Main paper written proofs with explicit model boundaries\\
Formal subset & C01, maintained NCT degree, and numerical C02; fresh compiler and kernel replay passed\\
New figures & 12 exact generated examples/formula plots, each with fixture provenance\\
Fresh source checks & Complete source \code{--all} verification and preserved historical outputs passed in the current replay\\
Rendering & This guide is compiled and visually checked separately from mathematical and kernel verification\\\bottomrule
\end{tabularx}

The source release, earlier companions, and archived evidence remain
unchanged. Publication corrections clarify evidence labels, the nonzero
monotone minterm domain, the $m=0$ scalar-summary boundary, and trace-variable
domains. No new conjecture or extension to unrestricted hardness is added.

\evidence{Figure and manuscript sources are included for review and reproduction.}

